\section{The Methods}

\begin{comment}
    

\begin{frame}{One Engine, Five Methods}
  \textbf{Shared core:} every method runs on the same incremental machinery ---
  per-move crossing deltas, spatial grids for candidate lookup, strict validity
  (no vertex on a foreign edge), and best-snapshot bookkeeping.
  \medskip
  \begin{center}
    \footnotesize
    \begin{tabular}{llll}
      \toprule
      \textbf{Method} & \textbf{Pipeline} & \textbf{Escapes local optima by} & \textbf{Binary} \\
      \midrule
      \texttt{sa}        & two-phase SA                    & temperature waves            & \texttt{sakgd} \\
      \texttt{sa-stress} & \textbf{stress init} $\to$ SA   & better start + waves         & \texttt{sfdp} + \texttt{sakgd} \\
      \texttt{staged}    & LNS (30\%) $\to$ SA (70\%)      & SA stage after greedy LNS    & \texttt{approach1} \\
      \texttt{staged-adaptive} & adaptive-LNS $\to$ SA     & growing neighbourhoods       & \texttt{approach1} \\
      \texttt{ils}       & SA rounds + random kicks        & perturbation between rounds  & \texttt{approach1} \\
      \bottomrule
    \end{tabular}
  \end{center}
  \medskip
  \textbf{Two-phase SA} (after the winning GD'25 entry):
  phase 1 minimises total crossings $\chi$ (untangles the drawing),
  phase 2 minimises the bottleneck $k$ with a dual fitness:
  $\Delta E = \Delta K_{\text{local}}$ if a touched edge's $k$ changes,
  else $\Delta\chi/\chi$ as a tiny tie-break.
\end{frame}

\end{comment}

\begin{frame}{Inside the SA Loop: Moves, Waves, Snapshots}
  \begin{columns}[T]
    \begin{column}{0.52\textwidth}
      \textbf{One move:}
      \begin{enumerate}
        \item pick a vertex --- crossing-weighted; in phase 2, biased to
              \emph{$k$-critical} endpoints,
        \item propose a position (Gaussian around current; radius shrinks
              with temperature),
        \item reject if occupied or vertex-on-edge,
        \item Metropolis: accept if $\Delta E\le 0$, else with
              probability $e^{-\Delta E/T}$.
      \end{enumerate}
      \medskip
      \textbf{Temperature waves:} inner loop cools $T$ per move;
      the outer loop restarts each wave from the best-so-far solution
      at a slightly lower starting temperature.
    \end{column}
    \begin{column}{0.45\textwidth}
      \textbf{Best-snapshot bookkeeping:}
      \begin{itemize}
        \item \texttt{saveBest()} --- whenever $(k,\chi)$ improves, copy the
              positions into \texttt{bestPos} (three fields, \emph{in RAM}).
        \item \texttt{restoreBest()} --- copy \texttt{bestPos} back and rebuild
              the tables: at wave ends, and once more before writing output.
      \end{itemize}
      \medskip
      \textbf{Stateless by design:} no files, no caches, nothing survives the
      process. Two runs on the same input are fully independent ---
      warm-starting across runs exists only as an explicit, off-by-default
      pipeline flag.
    \end{column}
  \end{columns}
\end{frame}

\begin{frame}{Method Shoot-Out: Same Seeds, Same Budgets}
  All five methods, all nine instances, one controlled sweep
  (4 workers, seeds 42--45, budgets 2/4/6\,min by size; best $k$ of 4 workers; all outputs validated):
  \begin{center}
    \footnotesize
    \begin{tabular}{lrrrrr}
      \toprule
      \textbf{Graph} & \texttt{sa} & \texttt{sa-stress} & \texttt{ils} & \texttt{staged} & \texttt{st.-adaptive} \\
      \midrule
      Automatic-1 & \textbf{9}  & \textbf{9}  & \textbf{9}  & \textbf{9}  & \textbf{9}  \\
      Automatic-2 & 4           & 4           & \textbf{3}  & \textbf{3}  & \textbf{3}  \\
      Automatic-3 & \textbf{35} & 37          & 37          & 37          & 37          \\
      Automatic-4 & 7           & \textbf{6}  & 7           & 7           & 7           \\
      Automatic-5 & 77          & \textbf{75} & 79          & 78          & 78          \\
      Automatic-6 & \textbf{669}& 737         & 787         & 754         & 765         \\
      Automatic-7 & 50          & \textbf{29} & 57          & 52          & 53          \\
      Automatic-8 & 47          & \textbf{8}  & 309         & 52          & 48          \\
      Automatic-9 & 57          & \textbf{11} & 53          & 59          & 56          \\
      \bottomrule
    \end{tabular}
  \end{center}
  \smallskip
  \begin{itemize}
    \item \texttt{sa-stress} wins or ties \textbf{6 of 9} --- and by $2$--$6\times$ where
          the canvas is tight (A7/A8/A9).
    \item Dense A6 is the exception: plain \texttt{sa} wins; stress init actively hurts there.
  \end{itemize}
\end{frame}

% --- slide 1 of 2: numbers ---
\begin{frame}{Staged-Adaptive \& ILS: The Numbers (A7--A9)}
  \texttt{sa-stress} vs.\ the two alternatives on the three hardest instances
  (4 workers, seeds 42--45, equal budgets; best $k$ of 4 workers):
  \bigskip
  \begin{center}
    \begin{tikzpicture}
      \begin{axis}[
        ybar, bar width=10pt,
        ymode=log,
        ylabel={best $k$ achieved (log scale, lower = better)},
        symbolic x coords={A7,A8,A9},
        xtick=data,
        xticklabels={A7 ($n{=}500$), A8 ($n{=}10\text{k}$), A9 ($n{=}3\text{k}$)},
        legend style={at={(0.5,-0.18)}, anchor=north, legend columns=3, font=\small},
        width=0.82\textwidth, height=5.8cm,
        ymajorgrids=true, grid style={dashed, gray!40},
        every axis/.append style={font=\small},
        nodes near coords,
        nodes near coords style={font=\scriptsize, inner sep=1pt},
        enlarge x limits=0.3,
      ]
        \addplot[fill=accent!80, draw=accent] coordinates {(A7,29) (A8,8)   (A9,11)};
        \addplot[fill=warm!80,   draw=warm]   coordinates {(A7,57) (A8,309) (A9,53)};
        \addplot[fill=muted!50,  draw=muted]  coordinates {(A7,53) (A8,48)  (A9,56)};
        \legend{\texttt{sa-stress} (winner),\quad\texttt{ils},\quad\texttt{staged-adaptive}}
      \end{axis}
    \end{tikzpicture}
  \end{center}
  \smallskip
  \textbf{A8 gap: $38\times$} --- \texttt{ils} reaches $k=309$, \texttt{sa-stress} reaches $k=8$.
\end{frame}

% --- slide 2 of 2: why ---
\begin{frame}{Staged-Adaptive \& ILS: Why They Lose}
  \begin{columns}[T]
    \begin{column}{0.48\textwidth}
      \textbf{ILS} (SA rounds $+$ random kicks)
      \medskip

      \begin{center}
      \begin{tikzpicture}[font=\small, node distance=0.7cm and 0.8cm,
          box/.style={draw, rounded corners=3pt, minimum width=2.2cm,
                      minimum height=0.55cm, align=center, inner sep=3pt}]
        \node[box, fill=teal!25]              (g)  {good layout};
        \node[box, fill=warm!35,  below=of g] (k)  {kicked layout};
        \node[box, fill=warm!15,  below=of k] (r)  {recovered (worse)};
        \draw[->, thick] (g) -- node[right=3pt]{\textit{random kick}} (k);
        \draw[->, thick] (k) -- node[right=3pt]{SA round} (r);

        \node[box, fill=teal!25, right=1.4cm of g]  (b1) {best-so-far};
        \node[box, fill=teal!40, below=1.4cm of b1] (b2) {refined best};
        \draw[->, thick, dashed, color=teal]
          (b1) -- node[right=3pt]{\shortstack{wave\\restart}} (b2);

        \node[above=4pt of g,  font=\footnotesize\bfseries, text=warm] {ILS};
        \node[above=4pt of b1, font=\footnotesize\bfseries, text=teal] {SA waves};
      \end{tikzpicture}
      \end{center}

      \smallskip
      On A8, untangling takes the whole budget.
      One kick resets that work.
      Waves restart \emph{from the best} --- no destruction.
    \end{column}
    \begin{column}{0.48\textwidth}
      \textbf{Staged-adaptive} (adaptive LNS $\to$ SA)
      \medskip
      \begin{itemize}
        \item LNS displaces many vertices at once, ignoring crossing geometry.
        \item The growing neighbourhood adds overhead, not quality.
        \item SA starts \emph{later} from a layout no better than stress init
              provides for free in seconds.
      \end{itemize}
      \vfill
      \textbf{Both dropped.}
      The complexity (LNS infrastructure, kick logic) adds nothing.
      \texttt{sa-stress} dominates on every large instance.
      The real lever was always the \emph{starting point}.
    \end{column}
  \end{columns}
\end{frame}
